\documentclass[10pt,a4paper]{amsart}
\usepackage[margin=19mm]{geometry}
\usepackage{amsmath,amssymb,mathtools}
\usepackage[scr=rsfs,cal=euler]{mathalfa}
\usepackage{xfrac}
\usepackage[unicode,hidelinks,bookmarks=false]{hyperref}
\newcommand{\ie}{i.e.\ }
\newcommand{\cf}{cf.\ }
\newcommand{\resp}{respectively\ }
\newcommand{\Subsection}{\subsection}
\providecommand{\nobreakdash}{\nobreak\hbox{-}}
\setlength{\emergencystretch}{2em}
\multlinegap0pt
\usepackage{mathtools}
\usepackage{amssymb}
\usepackage{bm}
\usepackage{xcolor}
\newtheorem{theorem}{Theorem}
\newtheorem{definition}{Definition}
\newtheorem{lemma}{Lemma}
\newcommand{\Div}{\mathrm{Div}}
\newcommand{\DivVal}{\mathrm{DivVal}}
\newcommand{\Fil}{\mathrm{Fil}}
\newcommand{\Nef}{\mathrm{Nef}}
\newcommand{\R}{\mathbb{R}}
\newcommand{\Spec}{\mathrm{Spec}}
\newcommand{\bR}{\mathbb{R}}
\newcommand{\cO}{\mathcal{O}}
\newcommand{\fX}{\mathfrak{X}}
\newcommand{\fa}{\mathfrak{a}}
\newcommand{\fm}{\mathfrak{m}}
\newcommand{\ord}{\mathrm{ord}}
\title{A note on b-divisors and filtrations on a local ring}
\author{Lu Qi}
\date{Source: arXiv:2604.07252v1 (CC BY 4.0)}
\begin{document}
\maketitle

\noindent\emph{Source and licence.} A note on b-divisors and filtrations on a local ring. Version arXiv:2604.07252v1 (CC BY 4.0), distributed by arXiv under the Creative Commons Attribution 4.0 International licence, http://creativecommons.org/licenses/by/4.0/, as recorded in the arXiv licence metadata for this exact version and shown on the abstract page https://arxiv.org/abs/2604.07252v1. Full licence notice: \texttt{licenses/arxiv-2604.07252-CC-BY-4.0.txt}.\par\smallskip

\noindent\textit{Editorial scope.} Counted unit: the article's theorem nef (source0.tex lines 384--411). The article's complete original proof of it is delivered with it (source0.tex lines 849--864, one \texttt{proof} environment of 16 lines, which the article places away from the statement). No mathematics is written, repaired or re-derived: the statement and the proof are the article's own lines, in the article's own notation, and no retained line is substituted or re-flowed.  Retained uncounted as the source-local context the proof needs: lines 332--334; lines 339--341; lines 622--629; lines 743--745; lines 817--819.  Nothing was omitted from any retained span, so the package carries no omission record.  No retained line contains a citation, so the wrapper carries no bibliography and no bibliography entry.  No retained line contains a figure environment, an image-inclusion command or drawing code, so no image is delivered and the package declares no asset dependency.  Editorial formatting only, in addition: an AMS \texttt{amsart} preamble that restates the article's own declarations and macros used by the retained lines, namely the environments \texttt{theorem}, \texttt{definition}, \texttt{lemma} on separate counters, together with the standard packages those lines need.  The retained environments print in this document as Theorem 1, Definition 1, Lemma 1 in order of appearance, whereas the article prints the same items with the numbering of its own full document.  The complete native source of the article as distributed in the arXiv e-print for this exact version is bundled unchanged as \texttt{sources/198081/source0.tex}.
\begin{equation}\label{eqn:filtr to b-div}
    Z_X(\fa_\bullet)\coloneqq -\sum_{P/X} \ord_P(\fa_\bullet)\cdot P,
\end{equation}

\begin{equation}\label{eqn:b-div to filtr}
    \fa(W)(U)\coloneqq \{f\in \cO_X(U)\mid Z(f)\le W\}
\end{equation}

\begin{definition}\label{defn:saturation}
    The \emph{saturation} $\widetilde{\fa}_\bullet$ of an $\fm$-filtration $\fa_\bullet$ is defined by
    \[
        \widetilde \fa_\lambda\coloneqq \{f\in \fm \mid v(f)\ge \lambda\cdot v(\fa_\bullet) \text{ for all } v\in\DivVal_{R,\fm} \}
    \]
    for each $\lambda\in \R_{>0}$.
    We say that $\fa_\bullet$ is \emph{saturated} if $\fa_\bullet=\widetilde\fa_\bullet$. 
\end{definition}

\begin{equation}\label{eqn:b-div to filtr local}
    \fa_\lambda(W):=\{f\in R\mid Z(f)\le \lambda W\}=\{v(f)\ge -\lambda \cdot v(W) \text{ for any } v\in\DivVal_{R,\fm}\}.
\end{equation}

\begin{lemma}\label{lem:bounded below gives filtration}
	For any $W\in\Div^+(\fX,x)$, the ideal $\fa_\lambda(W)=\fa(\lambda W)(X)$ is $\fm$-primary for any $\lambda\in\bR_{>0}$, and the collection $\{\fa_\lambda(W)\}_{\lambda\in\bR_{>0}}$ is an $\fm$-filtration, denoted by $\fa_\bullet(W)$. Moreover, if $W\in\Div^b(\fX,x)$, then $\fa_\bullet(W)\in\Fil^s_{R,\fm}$.
\end{lemma}

\begin{theorem}\label{thm:structure of fil and nef}
	Let $(R,\fm)$ be a normal, excellent Noetherian local domain. Write $x\in X$ for the closed point $\fm\in \Spec R$. Then the following statements hold.
	
	\begin{enumerate}
        \item The formula \eqref{eqn:filtr to b-div} gives an injective map
        \[
            Z_X(\bullet):\Fil^s_{R,\fm}\to\Div^b(\fX,x),
        \]
        which is continuous in the coefficientwise topology. %and the closure of the image is $\Nef_+(\fX,x)$, 
		
        \item The global sections $\fa_\lambda(W)\coloneqq \fa(\lambda W)(X)$ defined as in \eqref{eqn:b-div to filtr}
        for $\lambda\in\bR_{>0}$ give a surjective map
        \[
            \fa_\bullet(\bullet):\Div^b(\fX,x)\to\Fil^s_{R,\fm}.
        \]
		
        \item We have 
        \[
            \fa_\bullet(Z_X(\fa_\bullet))=\fa_\bullet
        \]
        for any $\fa_\bullet\in\Fil^s_{R,\fm}$, and 
        \[
            Z_X(\fa_\bullet(W))\le W
        \]
        for any $W\in \Div^b(\fX,x)$.
	\end{enumerate}
    
\end{theorem}

\begin{proof}[Proof of Theorem \ref{thm:structure of fil and nef}]
    By Lemma \ref{lem:bounded below gives filtration}, we know that $\fa_\lambda(W)\coloneqq\fa(\lambda W)(X)$ defines a map 
    \[
        \fa_\bullet(\bullet):\Div_b(\fX,x)\to \Fil^s_{R,\fm}.
    \]
    Moreover, in view of \eqref{eqn:filtr to b-div} and \eqref{eqn:b-div to filtr local}, the equality $\fa_\bullet=\fa_\bullet(Z_X(\fa_\bullet))$ is just a restatement of Definition \ref{defn:saturation} for saturated filtrations, from which the injectivity of $Z_X(\bullet)$ and the surjectivity of $\fa_\bullet(\bullet)$ follow.

    The continuity of $Z_X(\bullet)$ follows from the definition of the coefficientwise topology.

    Finally, for any $\lambda\in\bR_{>0}$ and $f\in\fa_\lambda(W)$, we know that $-v(f)\le \lambda v(W)$ for any $v\in\DivVal_{R,\fm}$. Hence $-v(\fa_\lambda(W))/\lambda\le v(W)$. 
    Letting $\lambda\to\infty$, we get the inequality 
    \[
        v(Z_X(\fa_\bullet(W)))=-v(\fa_\bullet(W))\le v(W)
    \]
    for any $v\in\DivVal_{R,\fm}$. The proof is finished.
\end{proof}

\end{document}
